% Supplementary material (v5). Same fallback note as main.tex: swap the two
% marked lines for \documentclass[sn-mathphys-num]{sn-jnl} when sn-jnl.cls is
% available. Tables in tables/*.tex are generated by
% src/synthclust/v5_tables.py; figures by src/synthclust/make_figures.py.
% Sections S7 and S8 are Appendices A and B of the v4 main text.
% Build variants: same two switches as main.tex (\SetReferee, \SetBlind);
% see joc_prep/README.md and ../build_submission.sh.
\ifdefined\SetReferee
  \documentclass[pdflatex,referee]{sn-jnl}
  \input{joc_prep/joc_referee}
\else
  \documentclass[pdflatex]{sn-jnl}
  % sn-jnl's \Authorfont sets \boldmath, which leaks into the abstract
  % (numbers and math printed bold); joc_referee.tex fixes the same leak.
  \makeatletter\g@addto@macro\abstractfont{\unboldmath}\makeatother
\fi
\input{joc_prep/joc_bib}
\providecommand\SetBlind{false}
\input{joc_prep/blind}
\pdfmapfile{=local.map}          %% build-machine workaround; delete elsewhere
\usepackage{graphicx}%
\usepackage{multirow}%
\usepackage{amsmath,amssymb,amsfonts}%
\usepackage{amsthm}%
\usepackage{mathrsfs}%
\usepackage[title]{appendix}%
\usepackage{xcolor}%
\usepackage{textcomp}%
\usepackage{manyfoot}%
\usepackage{booktabs}%

\newcommand{\ind}{\mathbf{1}}
\newcommand{\bx}{\mathbf{x}}
\newcommand{\bmu}{\boldsymbol{\mu}}

\renewcommand{\thefigure}{S\arabic{figure}}
\renewcommand{\thetable}{S\arabic{table}}
\renewcommand{\thesection}{S\arabic{section}}
\renewcommand{\theequation}{S\arabic{equation}}

\begin{document}

\title{Supplementary Material\\
       \large Cluster Analysis of CART-Synthesised Data: A Factorial
       Simulation of $k$-Means and Ward's Linkage}
\author[1]{Oriol Farr\'es Vilar}
\author[2]{Jordi Cort\'es}
\author[3]{Xinnuo Chen}
\author[4]{Victoria L\'opez}
\author[2]{Daniel Fern\'andez}
\affil[1]{Department of Computer Science,
          Universitat Polit\`ecnica de Catalunya, Barcelona, Spain}
\affil[2]{Department of Statistics and Operations Research,
          Universitat Polit\`ecnica de Catalunya, Barcelona, Spain}
\affil[3]{Department of Statistical Science,
          University College London, London, United Kingdom}
\affil[4]{CUNEF Universidad, Madrid, Spain}
\abstract{Supporting figures and tables for the main article, the bootstrap
null by stratum, the clustering objectives and the computational cost of the
study. Unless stated otherwise, the number of groups is selected over the
fixed candidate range $\{2,\dots,8\}$, as in the main article.}
\maketitle

\section{Recovery of the number of groups by cluster count and separation}
\label{sec:heatmap}

Figure~\ref{figS:heatmap} presents the recovery rate of the true number of
groups as a function of the true cluster count $k$ and the separation
$\sigma$, for both algorithms. It complements panel (a) of Figure~3 of
the main article in a different layout, making the $k \times \sigma$
interaction directly legible: recovery collapses in the lower-left quadrant
(many groups, little separation) and is high in the upper-right.

\begin{figure}[!htbp]
\centering
\includegraphics{figures/figS_heatmap.pdf}
\caption{Recovery rate of the true number of groups (share of synthetic
  datasets with $\hat{k}_S = k$) by true cluster count and separation, for
  $k$-means (left) and Ward (right). Cell values are proportions over the
  60{,}000 evaluations of each cell; candidate range $\{2,\dots,8\}$.}
\label{figS:heatmap}
\end{figure}

\section{Agreement in the estimated number of groups between original and
         synthetic data}
\label{sec:khat}

Table~\ref{tabS:khat} cross-tabulates the number of groups estimated from the
original data against the number estimated from the corresponding synthetic
data, stratified by separation. Unlike the recovery rate, which scores
$\hat{k}$ against the construction-time ground truth, this table compares the
two estimates with each other and so removes the difficulty of recovering the
planted structure; it still contains the ordinary resampling instability of
$\hat{k}$, which Section~\ref{sec:null} quantifies.

The pattern confirms the expectation that at high separation the mass
concentrates on the diagonal: exact agreement rises from $68.4\%$ at
$\sigma = 0.1$ to $98.2\%$ at $\sigma = 10$ for $k$-means (Ward: $64.6\%$ to
$97.5\%$). In the overlapping regime the off-diagonal mass is not symmetric.
At $\sigma = 0.1$ the synthetic data yield a \emph{lower} estimate than the
original in $15.2\%$ of $k$-means comparisons and a higher one in $16.4\%$; the largest
single off-diagonal cell is the original data supporting $\hat{k} \ge 5$ while
the synthetic data support $\hat{k} = 2$ ($8.4\%$). That is, when the original
data contain weak, diffuse structure that the silhouette criterion resolves
into many small groups, CART synthesis tends to smooth it away into a single
coarse split. The tendency disappears once the groups are genuinely separated.

\begin{table}[!htbp]
\centering
\caption{Agreement in the estimated number of groups between original and
  synthetic data, by separation. Rows are $\hat{k}$ from the original data,
  columns $\hat{k}$ from the synthetic data; cells are percentages of all
  comparisons at that separation (each block sums to $100\%$). Estimates of
  five or more groups are pooled for display, but the agreement column is
  computed on unpooled values. Both data sources are searched over the fixed
  range $\hat{k} \in \{2,\dots,8\}$ used in the main article; under the capped
  range $\{2,\dots,\min(k+3,8)\}$ of Section~\ref{sec:range}, agreement is
  higher in the overlapping regime ($75.1\%$ and $67.6\%$ at $\sigma = 0.1$,
  $79.7\%$ and $75.9\%$ at $\sigma = 2$, for $k$-means and Ward respectively)
  and essentially unchanged at $\sigma \ge 6$, where the two conventions
  differ by at most $0.1$ percentage points. $n = 180{,}000$ comparisons per
  separation level.}
\label{tabS:khat}
\small
\input{supp_table.tex}
\end{table}

\section{Sensitivity to the candidate range for the number of groups}
\label{sec:range}

The main article selects $\hat{k}$ over the fixed range $\{2,\dots,8\}$. Earlier
work used $\{2,\dots,\min(k+3,8)\}$, a rule that ties the search range to the
planted number of groups and so passes information about the true $k$ to the
estimator. Table~\ref{tabS:range} repeats the count analyses under that capped
range, applied identically to original and synthetic data, and
Table~\ref{tabS:nullcapped} repeats the bootstrap null. Narrowing the range
raises both agreement measures in the overlapping regime, where the silhouette
criterion occasionally selects large values of $\hat{k}$, and leaves them
essentially unchanged at $\sigma \ge 6$. The qualitative conclusions are
unaffected: the two algorithms remain indistinguishable against the planted
number of groups, $k$-means agrees more often with the original estimate, the
advantage arises on Gamma data ($0.833$ versus $0.749$) and not on Normal data
($0.896$ versus $0.898$), and it is matched by the bootstrap null ($0.891$
versus $0.841$). Under the capped range the difference between the two
algorithms in the excess over the null is not significant ($-0.009$,
$p = 0.204$). All Wilcoxon statistics in this material and in the main
article treat scenarios with exactly equal means as ties (differences are
rounded to $10^{-12}$ before ranking, so that floating-point noise does not
break genuine ties).

\begin{table}[!htbp]
\centering
\caption{Recovery of the planted number of groups (RR) and agreement with the
  number estimated from the original data (AG), by separation and
  distribution family, under the fixed candidate range of the main article and
  under the capped range. Wilcoxon $W$ and $p$: paired signed-rank test of
  $k$-means against Ward across the 144 scenario-level means. Last row: overall
  bootstrap null $\mathrm{AG}_0$ ($B = 30$) under each range.}
\label{tabS:range}
\small
\input{tables/tabS_range.tex}
\end{table}

\begin{table}[!htbp]
\centering
\caption{Bootstrap null under the capped candidate range
  $\{2,\dots,\min(k+3,8)\}$: agreement after CART synthesis (AG), after
  bootstrap resampling ($\mathrm{AG}_0$, $B = 30$) and the excess
  $\mathrm{AG} - \mathrm{AG}_0$ (computed before rounding). Compare with
  Table~3 of the main article.}
\label{tabS:nullcapped}
\small
\input{tables/tabS_null_capped.tex}
\end{table}

\section{Bootstrap null and independent-pairs baseline}
\label{sec:null}

To judge whether a given level of agreement is high, it helps to know how
often a dataset that differs from the original only by resampling yields the
same estimate. The main article uses a bootstrap null: each of the 720 original
datasets is resampled $B = 30$ times with replacement, the scaler fitted on
the original data is applied unchanged, the same silhouette scan is run, and
$\mathrm{AG}_0$ is the share of resamples whose estimate equals $\hat{k}_R$.
Because proper synthesis in \texttt{synthpop} fits its trees to a bootstrap
sample of the original data, the resample is the synthesis procedure without
the tree step. Table~\ref{tabS:baseline} sets $\mathrm{AG}$ and
$\mathrm{AG}_0$ beside a second, coarser reference retained from the earlier
analysis: the proportion of the 10 pairs of independent original datasets in
a scenario that yield the same $\hat{k}_R$. The centroids are redrawn for
every original dataset (only the minimum pairwise centroid distance is fixed
at $\sigma$), so two such datasets differ in their centroid geometry as well
as in the sample. The pairs baseline is therefore not a second sample from the
same population and rests on only 10 pairs per scenario; it is reported as a
secondary reference. Both references agree that most of the disagreement in
the overlapping regime is present without synthesis, but only the bootstrap
null holds the geometry fixed.

\begin{table}[!htbp]
\centering
\caption{Synthetic-versus-original agreement (AG), the bootstrap null
  ($\mathrm{AG}_0$, $B = 30$ resamples per original dataset) and the
  independent-pairs baseline (Pairs: proportion of pairs of independent
  original datasets in the same scenario with equal $\hat{k}_R$, which also
  differ in centroid geometry). Means over scenarios; candidate range
  $\{2,\dots,8\}$.}
\label{tabS:baseline}
\small
\input{tables/tabS_baseline.tex}
\end{table}

\section{Geometric floor and scale-free metrics by stratum}
\label{sec:floor}

Tables~\ref{tabS:floor} and~\ref{tabS:scalefree} report, for every geometric metric, the observed
value after CART synthesis, the bootstrap floor (the same metric computed at
the planted $k$ between the partition of the original data and the partition
of a bootstrap resample, averaged over the $B = 30$ resamples) and the excess
of the observed value over the floor. MCD, MVD and $\Delta G$ are defined in
the main article; $\mathrm{MVD}/p$ divides MVD by the number of coordinates;
$\mathrm{MCD}_{\mathrm{sd}} = \mathrm{MCD}/s_w$, where
$s_w = \{\sum_i \lVert \bx_i - \bmu_{g_i} \rVert^2 / (Np)\}^{1/2}$ is the
pooled within-planted-group standard deviation of the original dataset in the
scaled coordinates; and
$\mathrm{RVD} = k^{-1}\sum_j |v_j^S/v_j^R - 1|$ is the relative variance
difference over the aligned clusters. The observed scale-free metrics are
computed for all 720{,}000 original--synthetic pairs, and the observed MCD, MVD
and $\Delta G$ recomputed in the same pass coincide with the stored values.
Two patterns stand out. The excess of RVD grows with separation for both
algorithms, whereas absolute MVD declines, so the decline of MVD is a decline
of the floor and of the units. The observed values are lower for $k$-means in
every stratum, but the excess is not: overall it is larger for $k$-means for
MCD and $\mathrm{MCD}_{\mathrm{sd}}$ and equal for the other metrics, with no
consistent ordering between the algorithms across strata.

\begin{table}[!htbp]
\centering
\caption{Geometric floor by stratum for the metrics of the main article:
  observed value after CART synthesis (Obs.), bootstrap floor ($B = 30$) and
  excess (Obs.\ minus floor), averaged over the scenario means of each
  stratum, at the planted $k$. Strata: separation (36 scenarios each),
  distribution family (72 each) and dimensionality (48 each).}
\label{tabS:floor}
\footnotesize
\input{tables/tabS_floor.tex}
\end{table}

\begin{table}[!htbp]
\centering
\caption{Scale-free metrics by stratum: MVD per coordinate ($\mathrm{MVD}/p$),
  centroid distance in within-component standard deviations
  ($\mathrm{MCD}_{\mathrm{sd}}$) and relative variance difference (RVD).
  Columns and strata as in Table~\ref{tabS:floor}. Observed values use all
  720{,}000 original--synthetic pairs.}
\label{tabS:scalefree}
\footnotesize
\input{tables/tabS_scalefree.tex}
\end{table}

\section{Collapse cells}
\label{sec:collapse}

At low separation the silhouette scan often returns the smallest candidate,
$\hat{k} = 2$, on the original and the synthetic data alike, so that the two
estimates agree for a reason unrelated to the preservation of structure. We
call a scenario a \emph{collapse cell} if all five of its original datasets
give $\hat{k}_R = 2$ (with $n = 5$, the threshold of at least $90\%$ means all
five). Because $\hat{k} = 2$ is the correct answer when $k = 2$, the
substantive subset is $k \ge 3$ (96 scenarios, 480 original datasets).
Table~\ref{tabS:collapse} reports the share of these original datasets with
$\hat{k}_R = 2$, the number of collapse cells, the share of all 720{,}000
evaluations that fall inside them, the agreement inside and outside them, and,
over all $k$, the share of original datasets with $\hat{k}_R = 2$ at $p = 10$
and $\sigma \le 2$. All collapse cells lie at $\sigma \le 2$.

\begin{table}[!htbp]
\centering
\caption{Collapse cells for $k \ge 3$ under the fixed candidate range
  $\{2,\dots,8\}$ of the main article and under the capped range. Datasets $\hat{k}_R = 2$: share of
  the 480 original datasets with $k \ge 3$ for which the silhouette scan
  selects two groups. Collapse cells: number of the 96 scenarios with $k \ge 3$. Evaluations in
  cells: share of all 720{,}000 evaluations. AG inside and outside: agreement in the evaluations with
  $k \ge 3$ inside and outside the collapse cells. Last column: share of all
  original datasets (all $k$) with $\hat{k}_R = 2$ at $p = 10$ and
  $\sigma \le 2$.}
\label{tabS:collapse}
\small\setlength{\tabcolsep}{4pt}
\input{tables/tabS_collapse.tex}
\end{table}

\section{Clustering objectives}
\label{sec:objectives}

$k$-means partitions $N$ observations $\{\bx_i\}_{i=1}^N$, $\bx_i \in
\mathbb{R}^p$, into $k$ non-overlapping groups by minimising the within-cluster
sum of squares \citep{macqueen1967}. Writing $\mathcal{C} = \{C_1,\dots,C_k\}$
for the partition and $\bmu_j$ for the centroid of $C_j$,
\begin{equation}
  \arg\min_{\mathcal{C}} \sum_{j=1}^{k} \sum_{\bx_i \in C_j}
    \lVert \bx_i - \bmu_j \rVert^2,
  \qquad
  \bmu_j = \frac{1}{|C_j|}\sum_{\bx_i \in C_j} \bx_i .
  \label{eq:wcss}
\end{equation}
Lloyd's iteration alternates assignment and update steps until the partition
stabilises, converging to a local optimum.

Hierarchical clustering builds a dendrogram by successive merging
\citep{ward1963}. Each observation begins as a singleton; at every step the two
clusters whose merger least increases the total within-cluster error sum of
squares are combined. With
\begin{equation}
  \mathrm{ESS} = \sum_j \sum_{\bx_i \in C_j} \lVert \bx_i - \bmu_j \rVert^2,
  \qquad
  \Delta_{\text{Ward}}(C_a, C_b)
    = \frac{|C_a|\,|C_b|}{|C_a| + |C_b|}\,
      \lVert \bmu_a - \bmu_b \rVert^2 ,
  \label{eq:ward}
\end{equation}
a flat partition into $k$ groups is obtained by cutting the dendrogram at the
corresponding height. The two methods therefore minimise the same
within-cluster sum of squares, $k$-means by iterative reassignment and Ward's
linkage by greedy, irreversible merging.

\section{Computational complexity and cost}
\label{sec:complexity}

This section states what the study costs to run and why. Throughout, $N$ is the
sample size, $p$ the dimensionality, $k$ the number of groups, $I$ the number of
Lloyd iterations to convergence, $K = 7$ the number of candidate values of
$\hat{k}$, and $A = 2$ the number of algorithms.

\subsection{Cost of the components}

Table~\ref{tabS:complexity} collects the asymptotic cost of each component
together with the number of times the executed design invokes it. The
invocation counts follow from the design: $144$ scenarios $\times\ n = 5$
replicates $= 720$ units, each with $m = 1000$ synthetic draws, giving
$720{,}000$ original--synthetic pairs.

\begin{table}[!htbp]
\centering
\caption{Cost per invocation and invocation count. Exact minimisation of the
         $k$-means objective is NP-hard in general and remains so for $k = 2$
         in arbitrary dimension \citep{inaba1994}, so Lloyd's local-search
         heuristic \citep{lloyd1982} is used, restarted ten times. Sources:
         CART \citep{breiman1984}; Lloyd iteration
         \citep{macqueen1967,lloyd1982}; Ward linkage
         \citep{ward1963,murtagh2017}; silhouette \citep{rousseeuw1987};
         Hungarian alignment \citep{kuhn1955,jonker1987}.}
\label{tabS:complexity}
\small
\begin{tabular}{lp{4.2cm}r}
\toprule
Component & Cost per invocation & Invocations \\
\midrule
CART synthesis      & $O(p^{2} N \log N)$                      & $7.2 \times 10^{5}$ \\
Lloyd iteration     & $O(IkNp)$ per restart ($\times 10$)       & $5.77 \times 10^{6}$ \\
Ward linkage        & $O(N^{2} p)$ time, $O(N^{2})$ space      & $5.77 \times 10^{6}$ \\
Silhouette          & $O(N^{2} p)$                             & $1.15 \times 10^{7}$ \\
Hungarian alignment & $O(k^{3})$                               & $1.44 \times 10^{6}$ \\
\bottomrule
\end{tabular}
\end{table}

Two structural facts follow. First, the \emph{binding} term is $O(N^{2}p)$,
contributed by Ward's linkage and by the silhouette, and it is a property of the
sample size rather than of the design: doubling $N$ quadruples the work, whereas
doubling $m$ or the number of scenarios only doubles it. The $O(N^{2})$ working
set of the linkage is also what caps $N$ in practice, since each worker holds
its own distance structure; at $N = 1000$ this is roughly $8$~MB per worker and
irrelevant, at $N = 10^{4}$ it is $800$~MB and binding. Second, model selection
dominates the fit count: estimating $\hat{k}$ costs $K$ fits and $K$ silhouette
evaluations per algorithm per dataset, so the fidelity stage performs
$720 \times A(K{+}1) + 720{,}000 \times A(K{+}1) = 1.15 \times 10^{7}$
clustering fits rather than the $1.44 \times 10^{6}$ that fitting at a single
known $k$ would require. Where a cheaper formulation exists we did not pursue
it: Ward admits an $O(N^{2})$ nearest-neighbour-chain implementation
\citep{murtagh2017,mullner2011} and $k$-means can be seeded to bound its
expected error \citep{arthur2007}, so the figures below are conservative
relative to a tuned implementation.

One algorithmic choice does change the constant materially. The partition of
each original dataset is computed once per replicate and reused across that
replicate's $1000$ draws. Computing it inside the draw loop would multiply the
original-side work by $10^{3}$ and produce identical output, since the original
data do not vary within a replicate.

\subsection{Parallel execution}

The design is a set of independent units, so the computation is
embarrassingly parallel in the strict sense: no communication, no shared mutable
state, and a serial fraction confined to reading configuration and writing the
merged result. Amdahl's law is therefore not the operative constraint at this
scale; hardware contention is, and the implementation targets three specific
sources of it.

\emph{Thread oversubscription.} The clustering routines call into BLAS, which by
default spawns one thread per core \emph{inside each worker}. With $P$ worker
processes this creates $P \times C$ runnable threads on $C$ cores, and the
resulting context-switching and cache thrashing dominate the useful work. All
three stages therefore pin \texttt{OMP\_NUM\_THREADS},
\texttt{OPENBLAS\_NUM\_THREADS} and \texttt{MKL\_NUM\_THREADS} to $1$ before
importing any numerical library, leaving parallelism entirely to the process
pool. The effect is not marginal: on the reference machine, $18$ concurrent
worker processes, each loading one original dataset ($p = 10$, $k = 4$) and
running the silhouette scans of both algorithms, complete in $5.9$~s with
threads pinned and $52.3$~s without, a factor of $8.8$
(\texttt{scripts/checks\_v5/thread\_bench.py}; the times include interpreter
start-up and data loading).

\emph{Load imbalance.} Unit cost varies with $p$ and $k$, so a static partition
of the task list leaves workers idle at the tail. The synthesis stage instead
uses a work-stealing queue: tasks are serialised into a spool directory and each
worker claims one by \texttt{rename(2)} into a \emph{doing} directory. Rename
within a filesystem is atomic, so exactly one worker can win a given task and no
lock is required. Workers pull candidates in batches of $200$ rather than
listing the spool per task, which at $1.44 \times 10^{5}$ tasks reduces
directory-scan overhead from hours to seconds. The Python stages achieve the
same effect through a dynamically scheduled pool.

\emph{Restart cost.} Each unit checkpoints its own result on completion, so an
interrupted run resumes at the last completed unit rather than from the start.
At a stage cost of several hours this converts a failure from a total loss into
the loss of the units in flight.

\begin{figure}[!htbp]
\centering
\includegraphics{figures/figS_scaling.pdf}
\caption{Strong scaling of the model-selection workload on the reference
         machine, $72$ independent units, BLAS pinned to one thread per worker:
         (a) speed-up over serial execution, (b) parallel efficiency. The
         dashed vertical line marks the 12 physical cores. Scaling is
         near-linear to $8$ workers ($6.9\times$, $87\%$ efficiency) and the
         knee falls at $12$, the physical core count. Beyond it the
         additional workers occupy SMT siblings that share execution units with
         an already-saturated core, so throughput still rises but efficiency
         falls to $56\%$ at $18$ workers and $45\%$ at $24$. The production runs
         use $18$: past the efficiency knee, but chosen for wall time on an
         otherwise idle machine. At $56\%$ efficiency, the clustering stage
         costs several hours (Table~\ref{tabS:cost}) rather than the $2.0$~h
         that dividing its serial cost by $18$ would suggest.}
\label{figS:scaling}
\end{figure}

Figure~\ref{figS:scaling} shows the measured strong-scaling curve. The shape is
the expected one for a compute-bound, memory-light workload on a simultaneous
multithreading processor, and it makes the operating point explicit: $18$
workers buy a $10.0\times$ speed-up over serial execution, while $12$ would give
$8.8\times$ at much better efficiency. Neither choice is a contribution; the
point of measuring is that the default of one worker per hardware thread is not
obviously right, and here it is not.

\subsection{Measured cost}

\begin{table}[!htbp]
\centering
\caption{Wall-clock cost at $m = 1000$ on an AMD Ryzen~9 9900X (12~cores,
         24~threads, $46$~GiB) with $18$ workers. The evidence differs by row:
         the fidelity stage is instrumented; synthesis is the interval over
         which its $1.44 \times 10^{5}$ output files (one per scenario and draw,
         each holding that draw for all five replicates, $7.2 \times 10^{5}$
         synthetic datasets in total) were written, so it includes any
         idle time inside that interval; the clustering stage was not
         instrumented, so its entry is a range from two measurements, each on
         a uniform random sample of $3000$ of its $1.44 \times 10^{5}$ inputs,
         run at the same worker count and scaled to the full set ($3.6$~h in
         the original timing run; $4.8$~h in a re-measurement,
         \texttt{scripts/checks\_v5/time\_sample.py}). The
         bootstrap null of the main article (Section~\ref{sec:null}) took a
         further $363$~s on $16$ workers, and the scale-free metrics for all
         720{,}000 pairs $745$~s.}
\label{tabS:cost}
\small
\begin{tabular}{llrr}
\toprule
Stage & Evidence & Wall time & Unit cost \\
\midrule
Synthesis ($7.2 \times 10^{5}$ datasets)  & write interval & $1.24$ h & $6.2$ ms/dataset \\
Clustering ($7.2 \times 10^{5}$ pairs)    & sampled        & $3.6$--$4.8$ h & $18$--$24$ ms/pair \\
Fidelity ($1.15 \times 10^{7}$ fits)      & instrumented   & $4.44$ h & $1.4$ ms/fit \\
\midrule
Total & & $9.3$--$10.5$ h & \\
\bottomrule
\end{tabular}
\end{table}

The synthetic datasets occupy $21$~GB as uncompressed Parquet. Cost is linear in
$m$: the same design at $m = 100$ completes in about a tenth of the time, and
every marginal mean of Table~2 of the main article (under both candidate
ranges) and every overall and marginal mean of MCD, MVD and $\Delta G$ differs
from the $m = 1000$ value by at most $0.0052$, so the larger run buys precision
rather than a change in any conclusion.

\clearpage
\bibliography{refs}

\end{document}
